\section{Part IX: The Founder's Research Investment Operating System}

\subsection{65. The 6-question pre-approval memo}

Before any material research expenditure, answer: \textbf{What do we not know? Why does it matter? What observation would change our decision? Who can investigate competently? What is the maximum cost and time? What will remain useful if the hypothesis fails?} Add a legal/data and scientific-integrity gate. Fund a small experiment if downside is contained and the information gain matters.

For larger investments add market alternatives, full economic cost, finance source restrictions, publication/IP policy, risk owner, institutional capability gained and kill conditions.

\subsection{66. Weekly laboratory review}

The review is not a theater for showcasing success. Ask what was predicted, what was observed, which competing explanations remain, what data or tools failed, what needs replication, whether safety or consent was violated, and what next experiment is most informative. Separate exploratory ideas from confirmed claims. Document decisions and unresolved contradictions.

Track experiment throughput \textbf{alongside} quality, not in place of it. When an expensive experiment fails, investigate whether it was a valuable negative result or an avoidable protocol failure.

\subsection{67. Monthly financial and compliance review}

Reconcile project cash, restricted funds, unbilled costs, equipment purchases, staff allocation, sponsored commitments, IP obligations, pending approvals, and equipment maintenance. Review forecast-to-complete, not just spend-to-date. A project 50\% through its budget may be only 10\% through its necessary evidence.

Research administrators are productive when they prevent losses, protect participants, secure eligible reimbursement and make reporting truthful-not when they maximize form count.

\subsection{68. Quarterly portfolio review}

Rank projects by \textit{new evidence and future option quality} rather than by storytelling. Continue those whose next experiment has attractive information value and affordable cost. Repair poorly designed experiments with fixable issues. Pause projects missing legal access or key talent. Stop those where strong evidence makes the central hypothesis irrelevant, or where funding threatens existing obligations.

Preserve accumulated knowledge when closing: protocols, code rights, raw records under retention rules, lessons and appropriate publication.

\subsection{69. When to decline a sponsor}

Decline a deal if it requires research fraud, deceptive publication, illegal data use, unsafe experimentation, exploitation of trainees, undisclosed conflicts or unfinanceable matching commitments. Negotiate when terms unnecessarily prevent publication, monopolize unrelated future IP, or create indefinite exclusivity without sufficient value. Use independent legal and scientific review for consequential agreements.

The ability to refuse a sponsor is a form of intellectual sovereignty purchased with financial reserves.

\subsection{70. The research north star}

The strongest measure of research success is not funding won, papers counted, patent applications filed, or startup valuations announced. It is whether the institution can repeatedly formulate consequential questions, pursue them rigorously, correct itself, train others to do the same, and translate appropriate knowledge into useful work without compromising the integrity of inquiry.

The enterprise is becoming scientific when \textbf{its standards can reject the founder's favorite hypothesis} and survive that rejection.
